\documentclass[11pt]{bookest}

\usepackage[utf8]{inputenc}
\usepackage[T1]{fontenc}
\usepackage{lmodern}
\usepackage{booktabs}
\usepackage{natbib}

\title{Field Notes on Book Layout}
\author{Noor Fatima}
\date{\today}

\begin{document}

% bookest's default \maketitle does not switch to a bare page style itself
% (only its `article` option does), so the running header and page number
% would otherwise print on the title page. Suppressed here instead.
\pagestyle{empty}
\maketitle
\clearpage
\pagestyle{headings}

\tableofcontents

\chapter{Preparing a Manuscript}

\section{Structure Before Style}

A manuscript's chapter and section divisions should be settled before any
layout decision is made \citep{knuth1984texbook}. \texttt{bookest} extends
the standard \texttt{book} class with coloured titling, running headers and a
configurable cover, but the underlying sectioning commands are unchanged.

\section{Front and Back Matter}

Table~\ref{tab:matter} sketches the three regions of a typical manuscript and
what belongs in each.

\begin{table}[h]
  \centering
  \begin{tabular}{lp{7cm}}
    \toprule
    Region & Contents \\
    \midrule
    Front matter & Title page, table of contents, preface \\
    Body & Numbered chapters and sections \\
    Back matter & Appendices, bibliography, index \\
    \bottomrule
  \end{tabular}
  \caption{Regions of a manuscript and what each one holds.}
  \label{tab:matter}
\end{table}

\chapter{Working with the Extended Class}

\section{Colour and Headers}

\texttt{bookest} adds document-wide colour accents and header styling on top
of \texttt{book}, without changing how chapters are declared or how
cross-references resolve.

\section{Citing and Concluding}

The class does not alter bibliography handling either: a standard
\texttt{natbib} bibliography, as used here \citep{knuth1984texbook}, compiles
the same way it would under plain \texttt{book}.

\bibliographystyle{plainnat}
\bibliography{references}

\end{document}
